\newpage{}
\section{Representation}
\label{sec:representation}

The self-supervised representation module maps a normalized edge window 
($\tensor{x}_{u\,v}$) to a latent state ($\tensor{z}_{u\,v}$). 
\texttt{VICReg\_Rank4}%
\note{
  Here \texttt{Rank4} means that the implementation
  extends the original VICReg architecture by one additional tensor rank, so the
  encoder operates on rank-4 input tensors: $\{batch, channels, sequence, features\}$.
}
is an encoder trained with the Variance Invariance Covariance Regulation objective of
Bardes, Ponce, and LeCun~\cite{bardes2022vicreg}.\\


Internally, the output is produced in two stages. 
\[
\begin{aligned}
  \tensor{x}_{u\,v}
  &\in \R^{C \times H_x \times D_x} \\
  &\downarrow \\
  \Phi &\in \R^{H_x' \times D_e} \\
  &\downarrow \\
  \tensor{z}_{u\,v} &\in \R^{D_e}.
\end{aligned}
\]


First it encodes the input ($\tensor{x}$) into a time-indexed latent sequence of embeddings 
$\Phi \in \R^{H_x' \times D_e}$ 
and then it collapses into $\tensor{z}_{u\,v} \in \R^{D_e}$ with a temporal reduction. 
In more detail, the encoding considers the input and the mask%
\note{
  Invalid entries are not treated as market observations; they are zeroed or 
  excluded whenever the network forms time-level statistics. 
  In the current architecture the temporal length is preserved, 
  so we can make use of the mask again for the temporal reduction. 
}: 
\[
  \Phi(\tensor{x}_{u\,v};\,{}_{x}\tensor{m}_{u\,v}) :
  \R^{C \times H_x \times D_x}
  \times
  \{0,1\}^{C \times H_x}
  \longrightarrow
  \R^{H_x' \times D_e}.\\
\]
Where the slice $\Phi^{i} \in \R^{D_e}$ would be the latent state at sequence position $i$. \\

In our case $H_x' = H_x$ and we decided to apply a masked temporal reducer ($\operatorname{TemporalReducer}$) so that:
\[
  \operatorname{TemporalReducer}(\Phi;\,{}_{x}\tensor{m}_{u\,v}) :
  \R^{H_x \times D_e}
  \times
  \{0,1\}^{C \times H_x}
  \longrightarrow
  \R^{D_e}.
\]
A simple masked mean temporal reducer is:
\begin{equation}
\label{eq:masked-mean-reducer}
  \omega^i
  =
  \mathbf{1}
  \left\{
    \sum_{c=1}^{C} {}_{x}m_{u\,v}^{c\,i} > 0
  \right\},
  \qquad
  \tensor{z}_{u\,v}
  =
  \frac{
    \sum_{i=1}^{H_x}
    \omega^i\cdot
    \Phi(\tensor{x}_{u\,v};\,{}_{x}\tensor{m}_{u\,v})^{i}
  }{
    \sum_{i=1}^{H_x} \omega^i
  },
\end{equation}
Here $(\omega^i)$ marks whether input position $(i)$ is usable: it equals $1$
if any channel is valid there, and $(0)$ otherwise. Other reducers, such as a last-valid
state, attention pooling, or a learned readout, are also valid choices as long
as the output remains in $\R^{D_e}$. In our case, if no input position has $\omega^i=1$,
the reducer is undefined and the sample is not exported as a valid representation.\\

The final exported representation could be described as%
\note{
  For all the expressions batch notation is obtained by stacking these 
  sample-level objects along a leading sample axis.
}: 
\[
  \tensor{z}_{u\,v}
  =
  \operatorname{TemporalReducer}
  \left(
    \Phi(\tensor{x}_{u\,v};\,{}_{x}\tensor{m}_{u\,v});\,
    {}_{x}\tensor{m}_{u\,v}
  \right)
  \in
  \R^{D_e}.
\]

The training data ranges
over the full family $\{\tensor{x}_{u\,v} : (u,v)\in\mathcal{E}\}$, so the
encoder learns and can be used across all instruments. 
\textit{All directed edges share the same representation map.}












\subsection{Representation Network}
\label{subsec:network-structure}

Let $E_c$ denote the internal channel-expansion dimension, $F$ the fused channel
dimension, and $H_{\Phi}$ the hidden width of the temporal convolution stack.  The
network has the following shape flow.

First, invalid channel-time entries are masked:
\[
  \widetilde{x}_{u\,v}^{c\,i\,d}
  =
  {}_{x}m_{u\,v}^{c\,i}\cdot x_{u\,v}^{c\,i\,d}. 
  \in \R
\]
A depthwise convolution over the time axis extracts local channel-wise
patterns, followed by a ReLU activation%
\note{
  $\operatorname{DWConv}_{3\times 1}$ is a grouped convolution with one group
  per input channel.  Its kernel spans three time positions and one feature
  coordinate, so it filters each channel independently; channel mixing is
  deferred to the later fusion layer.
}$^{,}$%
\note{
  $\operatorname{ReLU}$ is applied coordinatewise.  For a scalar input $a$,
  $\operatorname{ReLU}(a)=\max(a,0)$.
}%
:
\[
  {}_{\Phi}\mat{U}^{(1)}
  =
  \operatorname{ReLU}
  \left(
    \operatorname{DWConv}_{3\times 1}(\widetilde{\tensor{x}})
  \right)
  \in
  \R^{C \times H_x \times D_x},
\]

Next, a grouped $1\times D_x$ convolution projects each channel's feature vector
into an $E_c$-dimensional latent channel space:
\[
  {}_{\Phi}\mat{U}^{(2)}
  \in
  \R^{C \times H_x \times E_c}.
\]
A learned channel identity vector $\vect{\kappa}_c \in \R^{E_c}$%
\note{
  The channel identity vector identifies the channel, not the edge or
  instrument.  The same $\vect{\kappa}_c$ is shared across all directed edges
  for channel $c$.
}
is added:
\[
  {}_{\Phi}\mat{U}^{(3),c\,i}
  =
  {}_{\Phi}\mat{U}^{(2),c\,i}
  +
  \vect{\kappa}_c.
\]
This gives the network an explicit way to distinguish, for example, the
\field{1m}, \field{1h}, and \field{1d} channels even when their normalized
feature values are similar. 
The tensor is then passed through the implementation's temporal transformer%
\note{
  This is a one-dimensional spatial-transformer-style module in the sense of
  Jaderberg et al.~\cite{jaderberg2015spatial}: it predicts sampling offsets
  from the current sequence, builds a differentiable grid along the time axis,
  samples the sequence, and then remasks invalid positions.  In the
  implementation, the channel and expanded-feature coordinates are temporarily
  flattened before sampling and reshaped back afterward.  It is not the
  attention Transformer architecture.
}%
:
\[
  {}_{\Phi}\mat{U}^{(4)}
  =
  \operatorname{TemporalTransformer1D}
  \left(
    {}_{\Phi}\mat{U}^{(3)};\,
    {}_{x}\tensor{m}_{u\,v}
  \right)
  \in
  \R^{C \times H_x \times E_c}.
\]
This block preserves the temporal length $H_x$ but learns a time-axis warp
before channel fusion.
And then the channel dimension is fused with a $1\times 1$ convolution%
\note{
  $\operatorname{Conv}_{1\times 1}^{C\rightarrow F}$ is a pointwise convolution
  applied at each time and expanded-feature location.  It mixes the $C$
  channel axis into $F$ fused channels without changing $H_x$ or $E_c$.
}%
:
\[
  {}_{\Phi}\mat{U}^{(5)}
  =
  \operatorname{ReLU}
  \left(
    \operatorname{Conv}_{1\times 1}^{C\rightarrow F}
    ({}_{\Phi}\mat{U}^{(4)})
  \right)
  \in
  \R^{F \times H_x \times E_c}.
\]
Then a linear projection maps the last coordinate from $E_c$ to $H_{\Phi}$:
\[
  {}_{\Phi}\mat{U}^{(6)}
  \in
  \R^{F \times H_x \times H_{\Phi}}.
\]
The $F$ and $H_{\Phi}$ coordinates are flattened into convolution channels,
giving
\[
  {}_{\Phi}\mat{U}^{(7)}
  \in
  \R^{(F H_{\Phi}) \times H_x}.
\]
Finally, a residual dilated one-dimensional convolution stack maps this sequence
to the representation width%
\note{
  $\operatorname{DilatedConv}$ is a stack of one-dimensional convolutional
  residual blocks along the time axis.  The dilation grows across the stack, so
  later blocks see a wider temporal context without shortening the sequence.
  The mask is reapplied through the stack so invalid time positions do not
  contribute as observations.
}%
:
\[
  \Phi(\tensor{x}_{u\,v};\,{}_{x}\tensor{m}_{u\,v})
  =
  \operatorname{DilatedConv}
  ({}_{\Phi}\mat{U}^{(7)};\,{}_{x}\tensor{m}_{u\,v})^{\top}
  \in
  \R^{H_x \times D_e}.
\]
The transpose notation only indicates that the final tensor is stored with time
as the middle axis and representation width as the last axis. At the end we can 
apply the temporal reduction in equation~\eqref{eq:masked-mean-reducer}.
























\subsection{VICReg training objective}
\label{subsec:vicreg-training-objective}

During self-supervised representation training, two \textbf{batched}
stochastic augmented views%
\note{
  For VICReg, the effective sample count is not only the batch size $B$.
  Because the loss is applied to valid time states after projection, the
  relevant count is $N_{\Omega}$, the number of valid $(b,i)$ pairs shared by
  the two augmented views.  Thus a moderate batch size, such as $B=64$, can
  still provide many training samples when each sequence contributes several
  valid time positions.  If $N_{\Omega}$ becomes too small relative to the
  projector width $D_p$, the variance and covariance estimates become unstable.
}
are generated from the same edge-labeled input batch.  We sample two
augmentation maps independently,
\[
  \operatorname{aug}_{(1)},\operatorname{aug}_{(2)}
  \overset{\mathrm{ind}}{\sim}
  \mathcal{A},
\]
and apply them to the same batch:
\[
  \begin{aligned}
    \left(
      \tensor{x}_{(1)},\,{}_{x}\tensor{m}_{(1)}
    \right)
    &=
    \operatorname{aug}_{(1)}
    \left(
      (\tensor{x}_{u\,v}^{b},\,{}_{x}\tensor{m}_{u\,v}^{b})_{b=1}^{B}
    \right),
    \\
    \left(
      \tensor{x}_{(2)},\,{}_{x}\tensor{m}_{(2)}
    \right)
    &=
    \operatorname{aug}_{(2)}
    \left(
      (\tensor{x}_{u\,v}^{b},\,{}_{x}\tensor{m}_{u\,v}^{b})_{b=1}^{B}
    \right).
  \end{aligned}
\]
The labels $(1)$ and $(2)$ mark the two augmented views.  They are semantic view
labels, not additional tensor axes.

The augmentation family $\mathcal{A}$ produces two alternative versions of the
same input $\tensor{x}$ and mask ${}_{x}\tensor{m}$%
\note{
  The augmentation map changes both the feature tensor and its validity mask.
  Dropped or masked positions are removed from the valid set rather than treated
  as observed zero-valued market data.
}.
It may perform distortions such as causal time warping, value jitter, time
masking, and point dropping%
\note{
  Be aware that the VICReg loss is computed on time positions that are valid in
  both augmented views.  This validity rule is stricter than the masked-mean
  reducer: a time position is used by VICReg only when all channels are valid in
  both augmented views.  This keeps the paired projected states comparable under
  the same channel support.
}.

VICReg loss requires and is applied after a projector MLP%
\note{
  This distinction matters because the VICReg loss is applied to the projection
  head output, while the forecast system later consumes the encoder
  representation.
}
with projector width $D_p$:
\[
  {}_{\zeta}\operatorname{Projector}
  :
  \R^{D_e}
  \rightarrow
  \R^{D_p}.
\]
The projector acts independently on each valid time state%
\note{
  The VICReg objective is applied to the time-indexed encoder states before the
  temporal reducer.  The reducer $\operatorname{TemporalReducer}$ defines the
  exported embedding used downstream, while VICReg learns from valid projected
  sequence states.  The projector sees one $D_e$-dimensional state at a time; it
  does not mix different time positions.  Coupling across valid $(b,i)$ states
  happens later through the mean, variance, and covariance terms of the VICReg
  loss.
}.
For each view $s\in\{1,2\}$,
\[
  {}_{\zeta}\vect{Y}_{(s)}^{b\,i}
  =
  {}_{\zeta}\operatorname{Projector}
  \left(
    \Phi
    \left(
      \tensor{x}_{(s)};
      {}_{x}\tensor{m}_{(s)}
    \right)^{b\,i}
  \right)
  \in
  \R^{D_p}.
\]
The projector is used for representation learning only.  It is not the object
consumed by the forecast head.

The valid training positions are the batch-time pairs that survive both
augmentations%
\note{
  If $N_{\Omega}\le 1$, the covariance estimate is not meaningful.  The
  implementation skips such batches rather than forcing the loss to be defined.
}:
\[
  \Omega
  =
  \left\{
    (b,i)
    :
    {}_{x}m_{(1)}^{b\,c\,i}
    \cdot
    {}_{x}m_{(2)}^{b\,c\,i}
    =
    1
    \text{ for all } c=1,\ldots,C
  \right\},
  \qquad
  N_{\Omega}=|\Omega|.
\]

For each valid pair $(b,i)\in\Omega$, the projector gives one vector from each
view.  Fix one row order over the valid set $\Omega$.  This row order introduces
a compact row index
\[
  n=1,\ldots,N_{\Omega}.
\]
If row $n$ corresponds to the valid pair $(b,i)\in\Omega$, define
\[
  {}_{\zeta}\vect{p}_{(s)}^{n}
  :=
  {}_{\zeta}\vect{Y}_{(s)}^{b\,i}
  \in
  \R^{D_p},
  \qquad
  s\in\{1,2\}.
\]
The same row order is used for both augmented views.  Thus
${}_{\zeta}\vect{p}_{(1)}^{n}$ and
${}_{\zeta}\vect{p}_{(2)}^{n}$ always come from the same original valid
batch-time pair $(b,i)$%
\note{
  The index $n$ is only the compact row index after masking.  The original
  valid positions are still the pairs $(b,i)\in\Omega$.  We use $n$ because the
  stacked projection matrix has rows $1,\ldots,N_{\Omega}$, while $\Omega$ may
  skip many original batch-time positions.
}.

For each view $s\in\{1,2\}$, define the stacked projection matrix
${}_{\zeta}\mat{P}_{(s)}$ by its entries:
\[
  {}_{\zeta}\mat{P}_{(s)}
  \in
  \R^{N_{\Omega}\times D_p},
  \qquad
  \left(
    {}_{\zeta}\mat{P}_{(s)}
  \right)^{n\,d}
  =
  \left(
    {}_{\zeta}\vect{p}_{(s)}^{n}
  \right)^{d},
  \quad
  n=1,\ldots,N_{\Omega},
  \quad
  d=1,\ldots,D_p.
\]

For each view $s\in\{1,2\}$, define the mean projected vector:
\[
  {}_{\zeta}\bar{\vect{p}}_{(s)}
  =
  \frac{1}{N_{\Omega}}
  \sum_{n=1}^{N_{\Omega}}
  {}_{\zeta}\vect{p}_{(s)}^{n}
  \in
  \R^{D_p}.
\]
The centered projection matrix is defined coordinatewise by
\[
  \left(
    {}_{\zeta}\widetilde{\mat{P}}_{(s)}
  \right)^{n\,d}
  =
  \left(
    {}_{\zeta}\vect{p}_{(s)}^{n}
  \right)^{d}
  -
  \left(
    {}_{\zeta}\bar{\vect{p}}_{(s)}
  \right)^{d},
  \qquad
  n=1,\ldots,N_{\Omega},
  \quad
  d=1,\ldots,D_p.
\]
Equivalently, this is the matrix obtained by subtracting the same mean vector
from every row of ${}_{\zeta}\mat{P}_{(s)}$.

The sample covariance of projected coordinates is
\[
  {}_{\zeta}\mat{C}_{(s)}
  \in
  \R^{D_p\times D_p},
\]
with entries
\[
  {}_{\zeta}C_{(s)}^{d\,d'}
  =
  \frac{1}{N_{\Omega}-1}
  \sum_{n=1}^{N_{\Omega}}
  \left(
    {}_{\zeta}\widetilde{\mat{P}}_{(s)}
  \right)^{n\,d}
  \left(
    {}_{\zeta}\widetilde{\mat{P}}_{(s)}
  \right)^{n\,d'},
  \qquad
  d,d'=1,\ldots,D_p.
\]

For each view $s\in\{1,2\}$ and projector coordinate
$d=1,\ldots,D_p$, define the coordinate column vector
\[
  {}_{\zeta}\vect{q}_{(s)}^{d}
  =
  \begin{bmatrix}
    \left({}_{\zeta}\vect{p}_{(s)}^{1}\right)^{d}\\
    \vdots\\
    \left({}_{\zeta}\vect{p}_{(s)}^{N_{\Omega}}\right)^{d}
  \end{bmatrix}
  \in
  \R^{N_{\Omega}}.
\]
Here $d$ indexes a projector coordinate, not a normalized kline feature
coordinate.

Using the same sample convention as the covariance matrix, define
\[
  \operatorname{Var}_{\Omega}
  \left(
    {}_{\zeta}\vect{q}_{(s)}^{d}
  \right)
  =
  \frac{1}{N_{\Omega}-1}
  \sum_{n=1}^{N_{\Omega}}
  \left[
    \left(
      {}_{\zeta}\vect{p}_{(s)}^{n}
    \right)^{d}
    -
    \left(
      {}_{\zeta}\bar{\vect{p}}_{(s)}
    \right)^{d}
  \right]^2.
\]
Equivalently,
\[
  \operatorname{Var}_{\Omega}
  \left(
    {}_{\zeta}\vect{q}_{(s)}^{d}
  \right)
  =
  {}_{\zeta}C_{(s)}^{d\,d}.
\]

Let $\gamma\in\Rpos$ be the VICReg variance floor%
\note{
  The variance floor is a target lower bound on projected-coordinate standard
  deviation:
  \[
    \sqrt{
      \operatorname{Var}_{\Omega}({}_{\zeta}\vect{q}_{(s)}^{d})
      +
      \eps
    }
    \ge
    \gamma.
  \]
  When this inequality holds for coordinate $d$, that coordinate contributes
  zero to the variance penalty.  The symbol $\gamma$ is local to the VICReg
  objective here; later edge-confidence weights such as $\gamma_e$ are different
  objects.
},
and let
\[
  \lambda_{\mathrm{inv}},
  \lambda_{\mathrm{var}},
  \lambda_{\mathrm{cov}}
  \in
  \Rnonneg
\]
be the term weights.  Here $[a]_+:=\max\{a,0\}$ denotes the positive part.  The
standard VICReg terms%
\note{
  The displayed equations state the clean VICReg objective.  The implementation
  uses numerically robust variants, including finite variance floors and
  covariance safeguards when the effective count $N_{\Omega}$ is small.
}
are
\[
  \mathcal{L}_{\mathrm{inv}}
  =
  \frac{1}{N_{\Omega}}
  \sum_{n=1}^{N_{\Omega}}
  \left\|
    {}_{\zeta}\vect{p}_{(1)}^{n}
    -
    {}_{\zeta}\vect{p}_{(2)}^{n}
  \right\|_2^2,
\]
\[
  \mathcal{L}_{\mathrm{var}}
  \left(
    {}_{\zeta}\mat{P}_{(s)}
  \right)
  =
  \frac{1}{D_p}
  \sum_{d=1}^{D_p}
  \left[
    \gamma
    -
    \sqrt{
      \operatorname{Var}_{\Omega}
      \left(
        {}_{\zeta}\vect{q}_{(s)}^{d}
      \right)
      +
      \eps
    }
  \right]_+,
  \qquad
  s\in\{1,2\},
\]
and
\[
  \mathcal{L}_{\mathrm{cov}}
  \left(
    {}_{\zeta}\mat{P}_{(s)}
  \right)
  =
  \frac{1}{D_p}
  \sum_{d=1}^{D_p}
  \sum_{\substack{d'=1\\d'\ne d}}^{D_p}
  \left[
    {}_{\zeta}C_{(s)}^{d\,d'}
  \right]^2,
  \qquad
  s\in\{1,2\}.
\]

The full self-supervised loss is
\[
  \mathcal{L}_{\mathrm{VICReg}}
  =
  \lambda_{\mathrm{inv}}\cdot\mathcal{L}_{\mathrm{inv}}
  +
  \lambda_{\mathrm{var}}\cdot
  \left(
    \mathcal{L}_{\mathrm{var}}({}_{\zeta}\mat{P}_{(1)})
    +
    \mathcal{L}_{\mathrm{var}}({}_{\zeta}\mat{P}_{(2)})
  \right)
  +
  \lambda_{\mathrm{cov}}\cdot
  \left(
    \mathcal{L}_{\mathrm{cov}}({}_{\zeta}\mat{P}_{(1)})
    +
    \mathcal{L}_{\mathrm{cov}}({}_{\zeta}\mat{P}_{(2)})
  \right).
\]
The invariance term makes the two views agree.  The variance term prevents
collapse to a constant representation.  The covariance term discourages
redundant projected coordinates.

After training, the forecast system uses the reduced embedding
$\tensor{z}_{u\,v}$.  The projector output
${}_{\zeta}\operatorname{Projector}(\cdot)$ is a training device and is not part
of the edge forecast distribution in Section~\ref{sec:edge-forecast-head}.